\documentclass[11pt]{article}

% ============================================================
% ENCODING AND LANGUAGE
% ============================================================
\usepackage[utf8]{inputenc}
\usepackage[T1]{fontenc}
\usepackage[french]{babel}

% ============================================================
% PAGE GEOMETRY
% ============================================================
\usepackage[margin=1in]{geometry}

% ============================================================
% MATHEMATICS
% ============================================================
\usepackage{amsmath}
\usepackage{amssymb}

% ============================================================
% TABLES
% ============================================================
\usepackage{booktabs}
\usepackage{longtable}
\usepackage{array}
\usepackage{tabularx}
\usepackage{adjustbox}
\usepackage{caption}
\usepackage{etoolbox}
\usepackage[table]{xcolor}

% Automatically center every table
\AtBeginEnvironment{table}{\centering}

% Center table captions
\captionsetup[table]{
    justification=centering,
    singlelinecheck=false
}

% More vertical spacing between rows
\renewcommand{\arraystretch}{1.25}

% Useful tabularx column type
\newcolumntype{Y}{>{\centering\arraybackslash}X}

% ------------------------------------------------------------
% RESULT COLORS
% ------------------------------------------------------------

\definecolor{winbg}{RGB}{218,242,221}
\definecolor{wintext}{RGB}{20,105,48}

\definecolor{tiebg}{RGB}{255,243,199}
\definecolor{tietext}{RGB}{130,92,0}

\definecolor{losebg}{RGB}{248,220,220}
\definecolor{losetext}{RGB}{155,35,35}

% Result labels
\newcommand{\WIN}{%
    \cellcolor{winbg}\textcolor{wintext}{\textbf{WIN}}%
}

\newcommand{\TIE}{%
    \cellcolor{tiebg}\textcolor{tietext}{\textbf{TIE}}%
}

\newcommand{\LOSE}{%
    \cellcolor{losebg}\textcolor{losetext}{\textbf{LOSE}}%
}

% Light horizontal separator between data rows
\newcommand{\rowrule}{\specialrule{0.30pt}{0pt}{0pt}}

% ============================================================
% NUMBERS
% ============================================================
\usepackage{siunitx}

\sisetup{
    round-mode=places,
    round-precision=3,
    scientific-notation=true
}

% ============================================================
% FIGURES AND LINKS
% ============================================================
\usepackage{graphicx}
\usepackage{hyperref}

% ============================================================
% DOCUMENT INFORMATION
% ============================================================
\title{
Définitions des métriques et résultats des benchmarks\\
\large
Schémas WENO/TENO augmentés par réseaux de neurones :
résumé numérique complet
}

\author{}
\date{\today}

% ============================================================
% DOCUMENT
% ============================================================

\begin{document}

\maketitle

\noindent
Tous les chiffres ont été régénérés en relançant directement les scripts
de comparaison du projet, sauf mention contraire
(Section~\ref{sec:strand2b}, où les modèles entraînés n'ont pas été
conservés et où les chiffres historiques sont cités depuis le journal de
résultats du projet).
Dans les tableaux comparatifs, les étiquettes
\textbf{WIN}, \textbf{TIE} et \textbf{LOSE}
sont toujours définies du point de vue du modèle neuronal :
WIN signifie que le modèle neuronal est meilleur que le schéma classique,
TIE qu'il est équivalent selon le critère considéré,
et LOSE qu'il est moins performant.

% ============================================================
% SECTION 1
% ============================================================

\section{Définitions des métriques}
\label{sec:metrics}

Soit $u_i$ la solution calculée par un schéma à la cellule $i$
(sur $N$ cellules), et $u_i^{\text{ref}}$ la solution réelle à la même
cellule.

Soit

\[
d_i = u_i-u_i^{\text{ref}}.
\]

Les normes utilisées sont :

\begin{align}
L_1
&=
\frac{1}{N}
\sum_{i=1}^{N}
|d_i|,
\\
L_2
&=
\sqrt{
\frac{1}{N}
\sum_{i=1}^{N}
d_i^2
},
\\
L_\infty
&=
\max_i |d_i|.
\end{align}

Pour les benchmarks du système d'Euler, elles sont calculées
indépendamment pour chaque variable conservée :

\begin{itemize}
    \item densité $\rho$ ;
    \item quantité de mouvement $\rho u$ ;
    \item énergie $E$.
\end{itemize}

Un résumé scalaire est ensuite calculé comme leur moyenne :

\[
\bar L_1
=
\frac{1}{3}
\left(
L_1^\rho
+
L_1^{\rho u}
+
L_1^E
\right).
\]


% ------------------------------------------------------------
\subsection{Variation totale}

\begin{equation}
\mathrm{TV}(u)
=
\sum_{i=1}^{N-1}
|u_{i+1}-u_i|.
\end{equation}

Comparée à la variation totale propre de la solution de référence,

\[
\mathrm{TV}(u) >
\mathrm{TV}(u^{\text{ref}})
\]

indique une oscillation parasite de type Gibbs introduite par le schéma,
absente de la vraie solution.


% ------------------------------------------------------------
\subsection{Erreur de conservation}

\begin{equation}
E_{\text{cons}}
=
\left|
\Delta x
\sum_{i=1}^{N}
\left(
u_i-u_i^{(0)}
\right)
\right|.
\end{equation}

où $u_i^{(0)}$ est la condition initiale.

Pour un schéma écrit sous forme de différence de flux sur un domaine
périodique ou fermé, cette quantité doit rester proche de la précision
machine quelle que soit la précision du schéma.

Une valeur non nulle significative indique donc principalement un
problème d'implémentation et non un simple déficit de précision numérique.


% ------------------------------------------------------------
\subsection{Dépassement / sous-dépassement
(overshoot / undershoot)}

\begin{equation}
\text{overshoot}
=
\max
\left(
0,
\max_i u_i-\max_i u_i^{\text{ref}},
\min_i u_i^{\text{ref}}-\min_i u_i
\right).
\end{equation}

Cette métrique mesure de combien la solution calculée dépasse le vrai
maximum ou passe sous le vrai minimum près d'une discontinuité.

Un schéma ayant un dépassement nul ne génère donc aucune oscillation
de Gibbs aux extrema, indépendamment de ses erreurs
$L_1$, $L_2$ ou $L_\infty$ ailleurs dans le domaine.


% ------------------------------------------------------------
\subsection{Positivité (Euler uniquement)}

\begin{equation}
\text{positif}
\iff
\min_i \rho_i > 0
\quad\wedge\quad
\min_i p_i > 0.
\end{equation}

Il s'agit d'un critère booléen et non d'une métrique de classement.

Un schéma produisant une densité ou une pression non positive a échoué
physiquement, indépendamment de la valeur de ses normes d'erreur.


% ------------------------------------------------------------
\subsection{Erreur de localisation du choc / contact}

Les localisations des structures sont trouvées à partir du gradient
de densité.

Le choc est identifié comme la position du plus fort

\[
\left|
\frac{\mathrm{d}\rho}{\mathrm{d}x}
\right|.
\]

Le contact est identifié comme le gradient le plus important en dehors
d'une petite fenêtre d'exclusion autour du choc.

L'erreur de localisation du choc est :

\begin{equation}
E_{\text{shock}}
=
\left|
x_{\text{shock}}^{\text{scheme}}
-
x_{\text{shock}}^{\text{ref}}
\right|.
\end{equation}


% ------------------------------------------------------------
\subsection{Critères réussite/échec utilisés pour déclarer un vainqueur}

\begin{align}
\texttt{beats\_teacher}
&\iff
L^{\text{NN}}
\le
L^{\text{classical}},
\\
\texttt{teacher\_fidelity}
&\iff
L^{\text{NN}}
\le
\phi
L^{\text{classical}},
\qquad
\phi = 1{,}02.
\end{align}

Le premier critère compare directement la précision du réseau
à celle du schéma classique.

Le second est un critère plus souple correspondant à l'idée de
« pas significativement pire ».

Le pourcentage d'amélioration ou de dégradation est calculé par :

\begin{equation}
X\%
=
\left(
\frac{L^{\text{NN}}}
     {L^{\text{classical}}}
-1
\right)
\times 100.
\end{equation}

Une valeur négative indique donc une erreur plus faible pour le modèle
neuronal, tandis qu'une valeur positive indique une erreur plus élevée.

Le critère \texttt{teacher\_fidelity} est uniquement utilisé dans
certains tests d'acceptation historiques de la section NN-TENO vs.\ TENO.

Il ne constitue pas une preuve de supériorité.


% ============================================================
% SECTION 2
% ============================================================

\section{NN-WENO vs.\ WENO}
\label{sec:strand1}

Reproduction de l'article d'origine.

Modèle :

\texttt{models/SeemsGood3.h5}.

Le réseau prédit directement la valeur reconstruite à l'interface,
selon l'approche de Stevens \& Colonius.

Les évaluations ont été réalisées via :

\begin{itemize}
    \item \texttt{Evaluate\_WENO5NN\_StepAdvection.py},
    \item \texttt{WENO5\_vs\_WENO5NN\_Burgers.py},
    \item \texttt{WENO5\_vs\_WENO5NN\_ShuOsher.py},
    \item \texttt{WENO5\_vs\_WENO5NN\_Sod.py}.
\end{itemize}


% ------------------------------------------------------------
\subsection{Advection d'un créneau (step advection)}

\begin{table}[htbp]

\begin{adjustbox}{max width=\textwidth,center}

\begin{tabular}{lrrc}

\toprule

\textbf{Métrique}
&
\textbf{WENO5}
&
\textbf{WENO5-NN}
&
\textbf{Résultat RN}
\\

\midrule

Variation totale max.\ (exact $=2{,}0$)
&
--
&
\num{2.00496}
&
\WIN
\\

\rowrule

Largeur du choc (cellules)
&
--
&
22 (sur 100)
&
\WIN
\\

\rowrule

Étendue de la solution
&
--
&
$[-0{,}00119,\ 1{,}00119]$
&
\WIN
\\

\bottomrule

\end{tabular}

\end{adjustbox}

\caption{
Advection d'un créneau, évaluation directe WENO5-NN
($n_x=100$, $n_t=451$).
\texttt{Evaluate\_WENO5NN\_StepAdvection.py}
ne rapporte que le côté RN directement.
Les chiffres comparatifs WENO5-vs.-RN
$L_1$/$L_2$/$L_\infty$/largeur-de-choc
($-28,7\%$/$-17,1\%$/$-5,2\%$/$-21,4\%$)
proviennent du journal de résultats déjà vérifié du projet
(\texttt{Report\_Results.md}, Partie 1.1).
}

\end{table}


% ------------------------------------------------------------
\subsection{Convergence sur l'équation de Burgers}

\begin{table}[htbp]

\begin{adjustbox}{max width=\textwidth,center}

\begin{tabular}{rrrrc}

\toprule

\textbf{Cellules}
&
\textbf{$\Delta x$}
&
\textbf{$E$ (WENO5-JS)}
&
\textbf{$E$ (WENO5-NN)}
&
\textbf{Résultat RN}
\\

\midrule

50
&
\num{4.0e-2}
&
\num{2.709410e-02}
&
\num{2.559760e-02}
&
\WIN
\\

\rowrule

100
&
\num{2.0e-2}
&
\num{1.681767e-02}
&
\num{1.606723e-02}
&
\WIN
\\

\rowrule

200
&
\num{1.0e-2}
&
\num{1.184928e-02}
&
\num{1.216283e-02}
&
\LOSE
\\

\rowrule

400
&
\num{5.0e-3}
&
\num{7.412588e-03}
&
\num{7.162685e-03}
&
\WIN
\\

\bottomrule

\end{tabular}

\end{adjustbox}

\caption{
Équation de Burgers avec une solution de référence à 1600 cellules.
$E$ représente l'erreur par rapport à la référence.
Le RN gagne sur trois des quatre résolutions testées ;
l'avantage n'est cependant pas monotone avec la résolution du maillage.
}

\end{table}


% ------------------------------------------------------------
\subsection{Tube à choc de Sod et Shu--Osher}
\label{sec:strand1-euler}

\begin{table}[htbp]

\begin{adjustbox}{max width=\textwidth,center}

\begin{tabular}{lrrrc}

\toprule

\textbf{Benchmark}
&
\textbf{WENO5}
&
\textbf{WENO5-NN}
&
\textbf{$\Delta$}
&
\textbf{Résultat RN}
\\

\midrule

Sod
&
\num{4.414673e-01}
&
\num{4.426384e-01}
&
$+0,265\%$
&
\LOSE
\\

\rowrule

Shu-Osher
&
\num{1.227778e+01}
&
\num{1.234543e+01}
&
$+0,551\%$
&
\LOSE
\\

\bottomrule

\end{tabular}

\end{adjustbox}

\caption{
$L_\infty$ par rapport à la solution réelle.
Pour Sod : solution exacte, 120 cellules.
Pour Shu-Osher : référence WENO5 à 3200 cellules,
300 cellules pour l'évaluation et $t=2{,}0$.
Les deux cas sont des quasi-égalités numériques,
mais restent légèrement défavorables au RN selon le critère strict.
}

\end{table}


% ============================================================
% SECTION 3
% ============================================================

\section{NN-TENO vs.\ TENO}
\label{sec:strand2}

Modèle NN-TENO5 entraîné par imitation.

Modèle :

\texttt{models/nn\_teno5.h5}

avec les seuils calibrés

\[
[0{,}66,\ 0{,}50,\ 0{,}69].
\]

Les évaluations ont été effectuées via :

\begin{itemize}
    \item \texttt{TENO5\_vs\_NNTENO5\_Scalar.py},
    \item \texttt{TENO5\_vs\_NNTENO5\_Euler.py}.
\end{itemize}

Pour les problèmes scalaires :

\[
N=200,
\qquad
\mathrm{CFL}=0{,}4.
\]

Pour les benchmarks Euler :

\[
N=120.
\]


% ------------------------------------------------------------
\subsection{Benchmarks scalaires}

\begin{table}[htbp]

\begin{adjustbox}{max width=\textwidth,center}

\begin{tabular}{lrrrc}

\toprule

\textbf{Benchmark}
&
\textbf{TENO5 $L_1$}
&
\textbf{NN-TENO5 $L_1$}
&
\textbf{$\Delta$}
&
\textbf{Résultat RN}
\\

\midrule

smooth\_sine\_advection
&
\num{3.327802e-07}
&
\num{3.327802e-07}
&
$0\%$
&
\TIE
\\

\rowrule

multifrequency\_advection
&
\num{2.433666e-05}
&
\num{2.433666e-05}
&
$0\%$
&
\TIE
\\

\rowrule

burgers\_smooth\_to\_shock
&
\num{1.328275e-04}
&
\num{1.328275e-04}
&
$0\%$
&
\TIE
\\

\rowrule

step\_advection
&
\num{1.331725e-02}
&
\num{1.340182e-02}
&
$+0,64\%$
&
\LOSE
\\

\bottomrule

\end{tabular}

\end{adjustbox}

\caption{
Benchmarks scalaires évalués avec le critère
\texttt{beats\_teacher}.
Trois des quatre cas sont des égalités exactes.
Le réseau reproduit donc la décision de TENO5 plutôt qu'une amélioration
indépendante.
}

\end{table}


% ------------------------------------------------------------
\subsection{Benchmarks Euler}

\begin{table}[htbp]

\begin{adjustbox}{max width=\textwidth,center}

\begin{tabular}{lrrrc}

\toprule

\textbf{Benchmark}
&
\textbf{TENO5 $\bar L_1$}
&
\textbf{NN-TENO5 $\bar L_1$}
&
\textbf{$\Delta$}
&
\textbf{Résultat RN}
\\

\midrule

Sod
&
\num{5.596243e-03}
&
\num{5.532622e-03}
&
$-1,14\%$
&
\WIN
\\

\rowrule

Shu-Osher
&
\num{7.344873e-02}
&
\num{7.307967e-02}
&
$-0,50\%$
&
\WIN
\\

\bottomrule

\end{tabular}

\end{adjustbox}

\caption{
Benchmarks Euler.
Les deux cas correspondent à des victoires réelles mais étroites du réseau.
Les marges restent faibles et ne constituent donc pas un avantage
systématique important.
}

\end{table}


% ------------------------------------------------------------
% DETAILED EULER TABLE
% ------------------------------------------------------------

\begin{table}[htbp]

\footnotesize

\begin{adjustbox}{max width=\textwidth,center}

\begin{tabular}{lrrrr}

\toprule

&
\multicolumn{2}{c}{\textbf{Sod}}
&
\multicolumn{2}{c}{\textbf{Shu-Osher}}
\\

\cmidrule(lr){2-3}
\cmidrule(lr){4-5}

\textbf{Métrique}
&
\textbf{TENO5}
&
\textbf{NN-TENO5}
&
\textbf{TENO5}
&
\textbf{NN-TENO5}
\\

\midrule

Densité $L_1$
&
\num{5.596243e-03}
&
\num{5.532622e-03}
&
\num{7.344873e-02}
&
\num{7.307967e-02}
\\

\rowrule

Quantité de mvt.\ $L_1$
&
\num{6.013741e-03}
&
\num{5.957161e-03}
&
\num{1.818382e-01}
&
\num{1.810088e-01}
\\

\rowrule

Énergie $L_1$
&
\num{1.356911e-02}
&
\num{1.341055e-02}
&
\num{3.667872e-01}
&
\num{3.622581e-01}
\\

\rowrule

Densité min.
&
\num{0.125}
&
\num{0.125}
&
\num{0.801257}
&
\num{0.801257}
\\

\rowrule

Pression min.
&
\num{0.1}
&
\num{0.1}
&
\num{1.0}
&
\num{0.998684}
\\

\rowrule

Positivité $\rho,p$
&
\multicolumn{2}{c}{\checkmark\ les deux}
&
\multicolumn{2}{c}{\checkmark\ les deux}
\\

\rowrule

Erreur loc.\ choc
&
\num{0.00833}
&
\num{0.00833}
&
0
&
0
\\

\rowrule

Erreur loc.\ contact
&
\num{0.00833}
&
\num{0.00833}
&
\num{0.0833}
&
\num{0.0833}
\\

\rowrule

Dépassement pression
&
\num{2.78e-17}
&
\num{2.78e-17}
&
\num{4.57e-04}
&
\num{1.33e-02}
\\

\bottomrule

\end{tabular}

\end{adjustbox}

\caption{
Décomposition complète par composante pour les benchmarks Euler.
Sur Shu-Osher, le dépassement de pression de NN-TENO5 est environ
29 fois plus important que celui de TENO5
(\num{1.33e-2} contre \num{4.57e-4}),
malgré la victoire du réseau sur la métrique $L_1$.
}

\end{table}


% ------------------------------------------------------------
\subsection{Tentative de rendre l'avantage réel, et non accidentel}
\label{sec:strand2b}

\emph{
Résultats historiques cités depuis \texttt{Results.md}, Phases 7--9.
Les modèles entraînés pour cette expérience n'ont pas été conservés et
n'ont donc pas pu être relancés pendant cette session.
}

Le modèle a été réentraîné sur des étiquettes oracle, avec un balayage du
régularisateur d'ancrage :

\[
\lambda_{\text{anchor}}
\in
\{1,5,20,100,500\}.
\]

L'évaluation est réalisée sur la même suite de six benchmarks utilisant
\texttt{beats\_teacher}.


\begin{table}[htbp]

\footnotesize

\begin{adjustbox}{max width=\textwidth,center}

\begin{tabular}{lccccccc}

\toprule

\textbf{$\lambda_{\text{anchor}}$}
&
\textbf{Sod}
&
\textbf{Shu-Osher}
&
\textbf{smooth}
&
\textbf{multifreq}
&
\textbf{Burgers}
&
\textbf{step}
&
\textbf{Total}
\\

\midrule

0 (non mélangé)
&
crash
&
n/a
&
\TIE
&
\TIE
&
\LOSE
&
explosion
&
catastrophique
\\

\rowrule

1
&
\LOSE
&
\LOSE
&
\TIE
&
\TIE
&
\WIN
&
\LOSE
&
3/6
\\

\rowrule

5
&
\LOSE
&
\LOSE
&
\TIE
&
\TIE
&
\WIN
&
\LOSE
&
3/6
\\

\rowrule

20
&
\LOSE
&
\LOSE
&
\TIE
&
\TIE
&
\TIE
&
\LOSE
&
2/6
\\

\rowrule

100
&
\WIN
&
\LOSE
&
\TIE
&
\TIE
&
\TIE
&
\LOSE
&
4/6
\\

\rowrule

500
&
\WIN
&
\LOSE
&
\TIE
&
\TIE
&
\TIE
&
\LOSE
&
4/6
\\

\midrule

\emph{référence sans oracle}
&
\WIN
&
\WIN
&
\TIE
&
\TIE
&
\TIE
&
\LOSE
&
\textbf{5/6}
\\

\bottomrule

\end{tabular}

\end{adjustbox}

\caption{
Aucune force d'ancrage ne rattrape ou ne dépasse la référence entraînée
par imitation pure, qui atteint 5/6.
Le réentraînement vers une cible locale théoriquement meilleure ne produit
donc pas automatiquement une meilleure trajectoire numérique.
}

\end{table}


% ============================================================
% SECTION 4
% ============================================================

\section{NN-TENO vs.\ TENO (Nouvelle approche recommandée par le professeur)}
\label{sec:strand-new}

Cette section couvre l'approche JAX par Unrolled Physics Loss, sur le cas
scalaire puis sur le système d'Euler.

% ------------------------------------------------------------
\subsection{Cas scalaire}
\label{sec:strand3}

Les résultats ont été générés via :

\texttt{ValidationMatrix\_ProjPDF.py}

avec :

\[
N=64.
\]

Le modèle utilisé est :

\texttt{run6\_learnedlambda}

un modèle \texttt{run3\_pretrained} réentraîné avec une tête de mélange
dédiée (mécanisme repris de la sous-section~\ref{sec:strand4}, Système
d'Euler, ci-dessous) et des exemples d'entraînement sur le flux non
linéaire de Burgers, absents de l'entraînement d'origine.


\begin{table}[htbp]

\begin{adjustbox}{max width=\textwidth,center}

\begin{tabular}{lrrrc}

\toprule

\textbf{Benchmark}
&
\textbf{TENO5}
&
\textbf{NN-TENO5 (JAX)}
&
\textbf{Rapport}
&
\textbf{Résultat RN}
\\

\midrule

Advection ($L_1$)
&
\num{3.481861e-02}
&
\num{3.381306e-02}
&
$0{,}9711\times$
&
\WIN
\\

\rowrule

Advection ($L_2$)
&
\num{9.895985e-02}
&
\num{9.552129e-02}
&
$0{,}9653\times$
&
\WIN
\\

\rowrule

Advection (TV, exact $=2$)
&
\num{2.0000}
&
\num{2.0000}
&
--
&
\TIE
\\

\rowrule

Burgers ($L_1$, avant-choc, hors distribution)
&
\num{2.069709e-04}
&
\num{1.787337e-04}
&
$0{,}8636\times$
&
\WIN
\\

\rowrule

Burgers ($L_2$, avant-choc, hors distribution)
&
\num{8.166970e-04}
&
\num{6.649269e-04}
&
$0{,}8142\times$
&
\WIN
\\

\bottomrule

\end{tabular}

\end{adjustbox}

\caption{
Évaluation à l'échelle canonique des benchmarks
($\text{final\_time}=1{,}0$ pour l'advection,
soit environ $8{,}5\times$ l'horizon utilisé pendant l'apprentissage à
$K=15$).
\textbf{NN-TENO5 (JAX) gagne sur les deux benchmarks}, y compris Burgers,
explicitement hors distribution (flux non linéaire jamais vu à
l'entraînement du tronc d'origine).
La tête de mélange dédiée démarre déjà gagnante avant tout entraînement
(mélange uniforme d'environ 4,7\% de la reconstruction du réseau à chaque
cellule) ; l'entraînement affiné améliore ensuite les deux résultats
davantage, sans jamais régresser.
}

\end{table}


% ------------------------------------------------------------
\subsection{Système d'Euler}
\label{sec:strand4}

Les résultats ont été générés via :

\texttt{TENO5euler\_vs\_NNTENO5euler\_JAX.py}.

Les paramètres sont :

\[
N=120,
\]

avec une référence à 320 cellules pour Shu-Osher.

Le meilleur modèle disponible utilise un tronc entraîné avec :

\begin{itemize}
    \item mélange de données lisses ;
    \item écrêtage du gradient ;
    \item tête de mélange dédiée entraînée séparément.
\end{itemize}

Le modèle correspondant est :

\texttt{..\_phase18\_learnedlambda.pkl}.


% ------------------------------------------------------------

\begin{table}[htbp]

\begin{adjustbox}{max width=\textwidth,center}

\begin{tabular}{lrrrc}

\toprule

\textbf{Benchmark}
&
\textbf{TENO5euler $\bar L_1$}
&
\textbf{NN-TENO5euler $\bar L_1$}
&
\textbf{Rapport}
&
\textbf{Résultat RN}
\\

\midrule

Sod
&
\num{8.258323e-03}
&
\num{8.287093e-03}
&
$1{,}0035\times$
&
\LOSE
\\

\rowrule

Shu-Osher
&
\num{2.131892e-01}
&
\num{2.092639e-01}
&
$0{,}9816\times$
&
\WIN
\\

\bottomrule

\end{tabular}

\end{adjustbox}

\caption{
Pour Sod, TENO5euler possède une erreur environ 0,35\% plus faible.
Pour Shu-Osher,
\textbf{NN-TENO5euler gagne de 1,84\%}.
Il s'agit de la première victoire nette contre un schéma classique
sur un benchmark canonique obtenue dans ce projet.
}

\end{table}


% ------------------------------------------------------------

\begin{table}[htbp]

\footnotesize

\begin{adjustbox}{max width=\textwidth,center}

\begin{tabular}{lrrrr}

\toprule

&
\multicolumn{2}{c}{\textbf{Sod}}
&
\multicolumn{2}{c}{\textbf{Shu-Osher}}
\\

\cmidrule(lr){2-3}
\cmidrule(lr){4-5}

\textbf{Métrique}
&
\textbf{TENO5euler}
&
\textbf{NN-TENO5euler}
&
\textbf{TENO5euler}
&
\textbf{NN-TENO5euler}
\\

\midrule

Densité $L_1$
&
\num{5.484873e-03}
&
\num{5.507869e-03}
&
\num{7.522187e-02}
&
\num{7.430708e-02}
\\

\rowrule

Quantité de mvt.\ $L_1$
&
\num{5.950426e-03}
&
\num{5.949739e-03}
&
\num{1.837471e-01}
&
\num{1.821342e-01}
\\

\rowrule

Énergie $L_1$
&
\num{1.333967e-02}
&
\num{1.340367e-02}
&
\num{3.805987e-01}
&
\num{3.713506e-01}
\\

\rowrule

Densité min.
&
\num{0.125}
&
\num{0.125}
&
\num{0.801257}
&
\num{0.800648}
\\

\rowrule

Pression min.
&
\num{0.1}
&
\num{0.1}
&
\num{1.0}
&
\num{1.0}
\\

\rowrule

Positivité $\rho,p$
&
\multicolumn{2}{c}{\checkmark\ les deux}
&
\multicolumn{2}{c}{\checkmark\ les deux}
\\

\rowrule

Erreur loc.\ choc
&
\num{0.00833}
&
\num{0.00833}
&
0
&
0
\\

\rowrule

Erreur loc.\ contact
&
\num{0.00833}
&
\num{0.00833}
&
\num{0.0833}
&
\num{0.0833}
\\

\rowrule

Erreur conservation (densité)
&
\num{1.33e-15}
&
\num{1.33e-15}
&
\num{3.79e-10}
&
\num{4.41e-10}
\\

\bottomrule

\end{tabular}

\end{adjustbox}

\caption{
Décomposition complète par composante.
NN-TENO5euler gagne sur deux des trois composantes de Shu-Osher
(quantité de mouvement et énergie).
La victoire agrégée sur $\bar L_1$ pour Shu-Osher provient donc
principalement de ces composantes plutôt que d'une amélioration uniforme
de toutes les variables conservées.
}

\end{table}


% ------------------------------------------------------------
\subsection{Progression : comment la victoire a été obtenue}

\begin{table}[htbp]

\begin{adjustbox}{max width=\textwidth,center}

\begin{tabular}{lrr}

\toprule

\textbf{Étape}
&
\textbf{Rapport Sod}
&
\textbf{Rapport Shu-Osher}
\\

\midrule

Avant (aucun modèle fini)
&
--
&
--
\\

\rowrule

Correctif de stabilité
(terme de perte mort + bug de réparation)
&
$1{,}29\times$
&
$1{,}24\times$
\\

\rowrule

+ données d'entraînement lisses
+ écrêtage du gradient
&
$1{,}16\times$
&
$1{,}23\times$
\\

\rowrule

+ reconstruction à mélange convexe
(littérature)
&
$1{,}145\times$
&
$1{,}065\times$
\\

\rowrule

+ tête de mélange dédiée,
entraînée séparément
&
$\mathbf{1{,}003\times}$
&
$\mathbf{0{,}982\times}$
\\

\bottomrule

\end{tabular}

\end{adjustbox}

\caption{
Progression obtenue grâce à plusieurs correctifs diagnostiqués
et vérifiés indépendamment.
Un rapport inférieur à $1$ signifie que le modèle neuronal possède
une erreur plus faible que TENO5euler.
}

\end{table}

Deux tentatives supplémentaires pour dépasser ce résultat sur Sod ont été
testées : un taux d'apprentissage 50 fois plus élevé pour la tête de mélange
(rapport Sod $1{,}0045\times$), puis une tête enrichie recevant directement
le stencil brut plutôt que les seules probabilités de trouble (rapport Sod
$1{,}0026\times$). Aucune ne dépasse une quasi-égalité. Un balayage direct
du coefficient de mélange (constante, sans apprentissage) confirme qu'il
s'agit d'une limite structurelle du tronc lui-même sur Sod, et non d'un
problème de calibration de la tête : l'erreur sur Sod augmente de façon
monotone avec tout coefficient de mélange testé (de $1{,}0000\times$ à
$\lambda{=}0$ jusqu'à $1{,}0472\times$ à $\lambda{=}0{,}3$), alors que
Shu-Osher s'améliore de façon monotone jusqu'à $\lambda \approx 0{,}2$.
Le modèle issu de la tête de mélange dédiée simple, présenté ci-dessus,
reste donc le meilleur disponible.


% ============================================================
% SUMMARY
% ============================================================

\section{Résumé}

\begin{table}[htbp]

\small

\begin{tabularx}{\textwidth}{
    @{}
    >{\raggedright\arraybackslash}p{0.25\textwidth}
    >{\raggedright\arraybackslash}p{0.25\textwidth}
    >{\centering\arraybackslash}p{0.11\textwidth}
    >{\raggedright\arraybackslash}X
    @{}
}

\toprule

\textbf{Volet}
&
\textbf{Verdict}
&
\textbf{Résultat}
&
\textbf{Observation principale}
\\

\midrule

1. WENO5-NN vs.\ WENO5
&
Dépend fortement du problème
&
\TIE
&
Victoire sur step advection et sur 3/4 résolutions de Burgers,
mais légère perte sur Sod et Shu-Osher.
\\

\rowrule

2. NN-TENO5 (imitation) vs.\ TENO5
&
Reproduit principalement le maître
&
\TIE
&
Trois égalités exactes sur les benchmarks scalaires,
deux petits gains sur Euler et une petite perte sur step advection.
\\

\rowrule

3. NN-TENO5 (JAX) vs.\ TENO5, scalaire
&
Victoire sur les deux benchmarks canoniques
&
\WIN
&
Gain de 3,5\% sur l'advection et de 18,6\% sur Burgers (hors distribution),
grâce à une tête de mélange dédiée et un entraînement incluant le flux de
Burgers.
\\

\rowrule

4. NN-TENO5euler (JAX) vs.\ TENO5euler
&
Première amélioration nette sur un benchmark canonique
&
\WIN
&
Victoire de $1,84\%$ sur Shu-Osher ;
Sod reste une quasi-égalité légèrement défavorable au RN.
\\

\bottomrule

\end{tabularx}

\caption{
Résumé global des quatre volets.
WIN signifie que le modèle neuronal démontre une amélioration,
TIE indique une performance globalement comparable,
et LOSE une performance inférieure au schéma classique.
}

\end{table}


Chaque variante RN, dans chaque volet, tourne aussi nettement plus lentement
par cellule que son équivalent classique :

\[
\sim 5\text{--}18\times
\]

pour les modèles basés sur Keras.

Les victoires en précision rapportées ici ne se traduisent donc actuellement
pas par un avantage pratique en temps de calcul.


% ============================================================
% END
% ============================================================

\end{document}